% Generated draft: requires mathematical and source review.
\documentclass[11pt]{article}
\usepackage[margin=25mm]{geometry}
\usepackage{fontspec}
\setmainfont{Latin Modern Roman}
\usepackage{amsmath,amssymb,mathtools,unicode-math}
\setmathfont{Latin Modern Math}
\usepackage{enumitem,needspace,xurl,hyperref}
\setlength{\parindent}{0pt}
\setlength{\parskip}{5pt}
\setlength{\emergencystretch}{3em}
\newcommand{\N}{\mathbb{N}}
\newcommand{\Z}{\mathbb{Z}}
\newcommand{\Q}{\mathbb{Q}}
\newcommand{\R}{\mathbb{R}}
\newcommand{\C}{\mathbb{C}}
\newcommand{\F}{\mathbb{F}}
\newcommand{\A}{\mathbb{A}}
\newcommand{\PP}{\mathbb{P}}
\newcommand{\E}{\mathbb{E}}
\newcommand{\eps}{\varepsilon}
\newcommand{\dd}{\,\mathrm{d}}
\begin{document}
\section*{UnsolvedMath ID: 30003400}
\textbf{Source problem code:} catalogue entry
\subsection{Definitions and mathematical statement}
For context, under the usual simple undirected graph convention, let \(G=(V,E)\), where \(E\subseteq [V]^2\). Its chromatic number \(\chi(G)\) is the least cardinal \(\kappa\) for which there exists a map \(c:V\to\kappa\) satisfying \(c(u)\ne c(v)\) whenever \(\{u,v\}\in E\). Its coloring number \(\operatorname{col}(G)\) is the least cardinal \(\kappa\) for which there exists a well-order \(\prec\) of \(V\) such that
\[
\left|\{u\in V:u\prec v\text{ and }\{u,v\}\in E\}\right|<\kappa
\qquad\text{for every }v\in V.
\]
These standard definitions provide context; the supplied record does not specify its graph convention or define its specialized two-parameter terminology.

The source asks: Is it consistent that there exists an infinite cardinal \(\theta\) and a \((\theta,\theta+2)\)-coloring graph? More generally, what is the relationship between compactness of chromatic number and compactness of coloring number?

The phrase ``\((\theta,\theta+2)\)-coloring graph'' is retained verbatim because the supplied record does not define it. In particular, it does not specify the graph size, the conditions imposed on the whole graph and its subgraphs, or the meaning of the two parameters. The notation \(\theta+2\) is also unresolved: for infinite cardinals, cardinal addition gives \(\theta+2=\theta\), whereas ordinal addition gives a different ordinal; interpreting it as the second successor cardinal \(\theta^{++}\) would require an unsupported correction. The background theory relative to which consistency is requested is not stated. The general compactness question does not specify cardinal thresholds, classes of subgraphs, bounds on either invariant, or the hypotheses under which implications are sought. Consequently, the supplied evidence supports the questions as written but does not determine a complete formal consistency assertion or a particular compactness implication.
\subsection{Short English statement}
The problem asks whether some infinite cardinal \(\theta\) can consistently support a graph described as a \((\theta,\theta+2)\)-coloring graph, and how compactness for chromatic number relates to compactness for coloring number. The specialized graph terminology, parameter notation, and intended compactness principles require clarification. The supplied dated curation labels the problem partially solved; this draft does not independently verify that assessment.
\subsection{Sources}
\begin{itemize}
\item Local UnsolvedMath registry: permanent ID \(30003400\), title ``Compactness of Chromatic and Coloring Numbers,'' source problem code \texttt{catalogue entry}, in the supplied Oberwolfach Reports Open Problems collection. The project identifies \texttt{data/unsolvedmath/problems.json} and \texttt{catalog/problems.json} as registry locations; this draft uses the supplied complete record. No individual proposer is supplied; the recorded proposed year is 2017.
\item Supplied source citation: \emph{Set Theory (2018)}, Oberwolfach Reports, pp.\ 527--589, DOI \texttt{10.4171/owr/2017/11}, \url{https://doi.org/10.4171/owr/2017/11}. The background describes Oberwolfach Reports as documenting mathematical workshops and their open-problem discussions. The citation's title year 2018 and the DOI's 2017 component are preserved without reconciliation.
\item Supplied literature source, labeled ``Oberwolfach report'': \url{https://ems.press/content/serial-article-files/46674}.
\item Supplied literature source, labeled ``Lambie-Hanson--Rinot'': \url{https://doi.org/10.1007/s00493-017-3741-6}. No title, theorem number, or precise supporting result is supplied.
\item Literature assessment dated exactly \texttt{2026-08-21}: the supplied curation labels the problem ``partially solved'' and reports substantial consistency and incompactness results distinguishing chromatic and coloring numbers and settling many cases, while describing the exact consistency question and general relationship as model-dependent and unresolved. This is supplied dated curation derived from the OWR record, not an independently verified literature assessment. The record expressly requests independent verification of time-sensitive status and cited claims.
\item The supplied statement status is \texttt{corrected\_verified}; its accompanying explanation says mathematical content was preserved while notation, typography, and sentence structure were repaired for publication. This records an editorial qualification and does not establish mathematical correctness or resolve the missing definitions.
\end{itemize}
\end{document}
